\documentclass{report}
\usepackage[utf8]{inputenc}
\usepackage[T1]{fontenc}
\usepackage{geometry}
\usepackage{graphicx}
\usepackage{float}
\usepackage{hyperref}
\usepackage[english,indonesian]{babel}

\title{Optimasi Pemrosesan Template LaTeX Menggunakan Pendekatan Berbasis AI}
\author{Academic Support Agent}
\date{\today}

\begin{document}
\maketitle

\begin{abstract}

\end{abstract}

\chapter{PENDAHULUAN}


\chapter{TINJAUAN PUSTAKA}


% Original content (commented out):
% \documentclass[a4paper, 11pt]{report}
\usepackage{setspace}
\setstretch{1,15}
\usepackage{newunicodechar}
\usepackage[french]{babel}
\usepackage{graphicx}
\graphicspath{{images/}}
\usepackage{fontspec}
\setmainfont{Verdana}
\usepackage{lipsum}
\renewcommand{\thesection}{\arabic{section}}
\usepackage{ulem}
\usepackage{geometry}
\geometry{hmargin=2.5cm,vmargin=3.5cm}

%%exemple titre : Un semestre à cracovie:
%% Entre orient et occident
%%De l'autre coté du rideau de fer
%%Associé à l'Est, veut se rapprocher de l'Ouest


\begin{document}

\begin{title}

\title{
    {\Huge Titre du devoir}\\
    \vspace{1 cm}
    {\LARGE Nom de l'élève Promo 20**}\\
    \vspace{2.5 cm}
    {\includegraphics{logo_efrei}}\\
    \vspace{2.5 cm}
    {\Large Nom du prof, matière/département}\\
    \vspace{1 cm}
    {\date{\Large Date du rendu}}}
\end{title}

\maketitle
\tableofcontents
\clearpage


\pagenumbering{arabic}

\section{Introduction}
\input{chapters/introduction}

\section{Développement}
\input{chapters/developpement/developpement}

\subsection{Partie 1}
\input{chapters/developpement/part1}

\subsection{Partie 2}
\input{chapters/developpement/part2}

\subsection{Partie 3}
\input{chapters/developpement/part3}

\section{Conclusion}
\input{chapters/conclusion}


\end{document}


\end{document}